\documentclass[12pt,onecolumn]{article}

\usepackage[utf8]{inputenc}
\usepackage[T1]{fontenc}
\usepackage[margin=1in]{geometry}
\usepackage{setspace}
\usepackage{amsmath,amssymb,amsthm}
\usepackage{tikz}
\usetikzlibrary{positioning,fit,backgrounds,arrows.meta}
\usepackage[colorlinks=true,linkcolor=blue,citecolor=blue,urlcolor=blue]{hyperref}
\usepackage{cleveref}
\usepackage{orcidlink}
\usepackage{fancyhdr}
\usepackage{xcolor}
\usepackage{graphicx}
\usepackage{mdframed}
\usepackage{xurl}

% ------------------------------------------------------------
% Basic formatting
% ------------------------------------------------------------

\setstretch{1.08}

\pagestyle{fancy}
\fancyhf{}
\setlength{\headheight}{15pt}
\lhead{Radio Reception as Local Registration}
\rhead{John C. W. McKinley}
\cfoot{\thepage}

\newcommand\blfootnote[1]{%
  \begingroup
  \renewcommand\thefootnote{}\footnote{#1}%
  \addtocounter{footnote}{-1}%
  \endgroup
}

\newcommand{\defterm}[2]{%
#1\textsuperscript{\textcolor{gray}{\hyperref[#2]{\scriptsize$\dagger$}}}%
}

% ------------------------------------------------------------
% Theorem environments
% ------------------------------------------------------------

\newtheorem{proposition}{Proposition}[section]
\newtheorem{definition}[proposition]{Definition}
\newtheorem{remark}[proposition]{Remark}

\crefname{proposition}{Proposition}{Propositions}
\Crefname{proposition}{Proposition}{Propositions}
\crefname{definition}{Definition}{Definitions}
\Crefname{definition}{Definition}{Definitions}
\crefname{remark}{Remark}{Remarks}
\Crefname{remark}{Remark}{Remarks}

% ------------------------------------------------------------
% Title
% ------------------------------------------------------------

\title{\textbf{Radio Reception as Local Registration:\\
A Device-Scale Instance of the Timeless Light Model}\\[0.4em]
\large\textit{Why Successful Propagation Calculations Do Not License Intermediate Life}}

\author{John C. W. McKinley\,\orcidlink{0009-0005-7097-5035}}
\date{September 2026}

\begin{document}

\maketitle

\blfootnote{\scriptsize
Working draft published at
\href{https://johncwmckinley.com}{johncwmckinley.com}.
This draft has not been deposited on Zenodo and may be revised.
}

% ------------------------------------------------------------
% Abstract
% ------------------------------------------------------------

\begin{abstract}

This paper gives a device-scale account of radio reception under the clarified architecture of the Timeless Light Model (TLM). The canonical photon result denies a persisting intermediate photon history between registered endpoints. The broader architecture distinguishes timeless lawful structure from relativistic registration: quantum mechanics describes lawful eligibility, while relativity governs the coherent spacetime organization of registered outcomes. Applied to radio, this means that reception is not fundamentally the interception of a temporally persisting signal-object that has crossed the intervening space. A transmitter-side event and a receiver-side event are lawfully related under Maxwellian and relativistic structure, but the received electrical outcome is realized locally in the receiver. The antenna, tuning stages, demodulator, and output circuitry register and transform local voltages and currents according to that lawful source--receiver relation. Standard electromagnetic propagation calculations, antenna theory, filtering, modulation, and demodulation remain unchanged. What is removed is only the additional ontology of a signal-object enduring through intermediate locations. Radio reception is therefore a clear everyday instance of the clarified architecture: lawful structure specifies the admissible relation; relativistic registration governs its spacetime appearance; and the receiver registers the outcome locally.

\end{abstract}

\noindent{\small\textcolor{gray}{
Technical terms marked with a superscript dagger
(\textsuperscript{\scriptsize$\dagger$})
link to their formal definitions in \Cref{sec:definitions}.
}}

\vspace{0.8em}

\begin{center}

\textit{The calculation spans the interval.}

\vspace{0.4em}

\textit{The registered events occur locally.}

\vspace{0.4em}

\textit{The calculation does not license an intermediate life.}

\end{center}

% ============================================================
\section{Introduction}
% ============================================================

Radio is commonly described in traveler language.

A transmitter emits a signal. The signal travels through space. The
signal reaches an antenna. The receiver catches it, extracts its
information, and produces sound, data, or another output.

That language is extraordinarily useful in engineering practice. It
compresses a large body of electromagnetic calculation into an intuitive
story.

The Timeless Light Model asks a narrower question:

\begin{quote}
What ontology is actually licensed by the successful calculation?
\end{quote}

The minimal photon-side canon of TLM begins with the null case. Along a
null worldline,

\[
d\tau = 0.
\]

No proper-time sequence accumulates for the photon. There is also no
photon rest frame from which an ordinary timelike persistence story can
be constructed. TLM therefore declines to add an internally lived
sequence of intermediate locations to the photon description
\cite{bedrock,retrocausality}.

This is the \defterm{refusal of transit}{def:refusal}.

The refusal of transit does not reject the null relation between source
and detector. It does not reject Maxwell's equations. It does not reject
retarded solutions, phase delay, diffraction, interference, attenuation,
antenna theory, or finite causal delay. It does not alter a radio
engineer's calculations.

It rejects one further step:

\begin{quote}
Because a mathematical solution assigns structured relations across
spacetime, there must therefore exist a persisting physical carrier whose
biography occupies those intermediate locations.
\end{quote}

That inference is not required by the formalism.

Radio provides an unusually clear application because the measurable
receiver-side event is plainly local. A voltage appears at the receiving
antenna. Currents change in the receiver's conductors. Filters, mixers,
demodulators, amplifiers, processors, and speakers undergo local state
changes.

Whatever the complete source--receiver relation may be, the electrical
quantity actually registered by the radio is the radio's own local
physical state.

The central claim of this paper is therefore:

\begin{mdframed}[
linewidth=0.8pt,
innertopmargin=10pt,
innerbottommargin=10pt,
innerrightmargin=14pt,
innerleftmargin=14pt,
backgroundcolor=gray!15,
leftmargin=0.33in,
rightmargin=0.33in
]
\normalsize\itshape

Standard electrodynamics lawfully relates transmitter-side and
receiver-side events. The receiver realizes its electrical outcome
locally. The success of the propagation formalism does not require a
persisting signal-object with an intermediate life between those
registrations.

\end{mdframed}

This is an interpretive restriction, not a revision of radio physics.

% ============================================================
\section{The Scope of the Claim}
% ============================================================

The argument is easiest to evaluate if its scope is kept narrow.

This paper does \emph{not} claim that radio reception is instantaneous.

It does \emph{not} claim that distance is irrelevant.

It does \emph{not} claim that shielding, geometry, antenna orientation,
frequency, phase, attenuation, reflection, refraction, diffraction, or
interference cease to matter.

It does \emph{not} claim that Maxwellian propagation calculations may be
discarded.

It does \emph{not} propose new radio equations.

It makes only the following ontological restriction.

\begin{proposition}[No-intermediate-life restriction]
\label{prop:nointermediatelife}

A successful formal description relating a transmitter-side event to a
receiver-side event does not by itself license a temporally persisting
signal-object whose identity occupies successive intermediate locations.

\end{proposition}

\begin{proof}

The formalism determines relations among sources, admissible
electromagnetic configurations, causal structure, detector conditions,
and measurable outcomes.

To infer from those relations the additional existence of one persisting
carrier requires an ontological premise not supplied merely by the
mathematical success of the relation.

In the photon case, null proper time and the absence of a rest frame
supply positive reason not to import the ordinary timelike persistence
template.

Therefore the successful source--receiver calculation remains intact,
while the additional carrier biography is withheld.

\end{proof}

\begin{remark}

The proposition does not say that the mathematical description of the
intermediate region is useless, fictitious, or dispensable.

It says that mathematical description and ontological biography are not
the same assertion.

\end{remark}

% ============================================================
\section{Lawful Relation and Registered Outcome}
% ============================================================

The clarified TLM architecture distinguishes lawful structure from the
registered physical outcomes governed by that structure
\cite{unified,generaltimelessness}.

The laws are not travelers.

Boundary conditions are not travelers.

A Maxwellian solution is not made into a traveler merely because its
variables possess spacetime dependence.

A wave equation may correctly determine what is registered at a later
location without thereby requiring the mathematical wave description to
be reified as one persisting thing living through each point of the
interval.

The transmitter contributes to the lawful physical conditions under
which later receiver-side registrations are possible.

The receiver then realizes a local physical result.

Schematically,

\[
\boxed{\text{transmitter-side registration}}
\]

\[
\Downarrow
\]

\[
\boxed{
\begin{array}{c}
\text{lawful source--receiver relation}\\
\text{described by Maxwellian and relativistic structure}
\end{array}}
\]

\[
\Downarrow
\]

\[
\boxed{\text{receiver-side registration}}.
\]

The middle term is indispensable.

But indispensability of the relation does not imply the existence of an
additional enduring object in the relation.

\begin{definition}[Lawful source--receiver relation]
\label{def:lawfulrelation}

A lawful source--receiver relation is the complete physical relation
specified by the governing laws, source conditions, boundary conditions,
geometry, and admissibility structure linking source-side conditions to
possible detector-side registrations.

\end{definition}

\begin{definition}[Intermediate life]
\label{def:intermediatelife}

Intermediate life is the attributed persistence of one physical
signal-object, photon-object, wave-object, or carrier-object through a
sequence of intermediate spacetime locations between registered
endpoints.

\end{definition}

\begin{definition}[Refusal of transit]
\label{def:refusal}

The refusal of transit is the TLM restriction against inferring an
intermediate life merely from the successful formal relation between
registered endpoints.

\end{definition}

The refusal is therefore not a refusal of relation.

It is a refusal to turn relation into biography.

% ============================================================
\section{What the Receiver Actually Registers}
% ============================================================

At the device level, radio reception is local.

A receiving antenna develops measurable electrical conditions at its own
terminals.

Its conductors contain charges belonging to the receiving apparatus.

Those charges participate in local electrical state changes.

The receiver's front end then amplifies, filters, mixes, or otherwise
transforms locally available electrical quantities.

A demodulator produces another local electrical state.

An output stage produces still further local changes: speaker motion,
stored data, display output, or control action.

\begin{proposition}[Locality of the received electrical outcome]
\label{prop:localoutcome}

The quantity called the received radio signal is realized as local
voltages, currents, and state changes in the receiving apparatus.

\end{proposition}

\begin{proof}

Every directly processed quantity in the receiver is instantiated in the
receiver's own physical components.

The antenna voltage is the voltage at the antenna.

The receiver current is the current in the receiver.

The demodulated voltage is produced by the receiver's own circuitry.

The speaker motion is produced by the receiver's own output stage.

The distant source enters through the lawful physical relation governing
which local responses are realized, but the measured electrical
quantities themselves are local properties of the receiving system.

\end{proof}

\begin{definition}[Local receiver realization]
\label{def:localrealization}

Local receiver realization is the physical registration within the
receiver of the electrical state changes constituting reception.

\end{definition}

This point does not depend upon any older TLM account of matter as a
special ``cadence,'' recurrence rhythm, or continuously self-generated
signal.

The claim is simpler.

The receiver is a physical system.

Its registered electrical state is local.

That local state is lawfully related to remote source conditions.

Nothing further is required for the present argument.

% ============================================================
\section{Definitions}
\label{sec:definitions}
% ============================================================

\begin{definition}[Lawful structure]
\label{def:lawfulstructure}

Lawful structure is the set of physical laws, constraints, symmetries,
conservation requirements, coupling rules, boundary conditions, and
mathematical relations determining which physical relations and outcomes
are admissible.

Lawful structure is not itself a temporal traveler.

\end{definition}

\begin{definition}[Lawful qualification]
\label{def:qualification}

Lawful qualification is the admissibility structure under which a
physical outcome or relation is permitted by the governing laws,
including conservation conditions, coupling conditions, boundary
conditions, and, where applicable, quantum amplitude and probability
structure.

\end{definition}

\begin{definition}[Registration]
\label{def:registration}

Registration is a physically realized outcome possessing definite
participation in spacetime description, including localization, temporal
ordering, causal relation, and the possibility of physical record or
receipt.

\end{definition}

\begin{definition}[Relativistic registration]
\label{def:relregistration}

Relativistic registration is registration organized according to the
causal, geometric, and temporal structure of special or general
relativity.

\end{definition}

\begin{definition}[Transit ontology of radio]
\label{def:transit}

Transit ontology of radio is the additional interpretation according to
which successful radio propagation requires one persisting
signal-bearing entity to occupy successive intermediate locations between
transmitter and receiver.

\end{definition}

\begin{definition}[Received signal]
\label{def:receivedsignal}

The received signal is the structured local electrical outcome
registered and processed by the receiving apparatus under the lawful
conditions relating it to the source.

\end{definition}

% ============================================================
\section{The Engineering Sequence}
% ============================================================

Nothing in the refusal of transit changes the engineering sequence.

\subsection*{Source}

A transmitter undergoes local time-dependent electrical changes.

Those changes serve as source conditions in the standard electromagnetic
formalism.

Modulation schemes alter the relevant source-side structure in precisely
the way described by conventional engineering.

\subsection*{Source--receiver calculation}

Maxwell's equations and the relevant boundary conditions determine the
lawful electromagnetic relation among the source, environment, and
possible receiving locations.

The resulting mathematics contains spatial and temporal dependence.

It supports calculations of radiation pattern, phase, attenuation,
reflection, diffraction, polarization, delay, and interference.

TLM retains all of these calculations.

What TLM does not add is:

\begin{quote}
There must therefore exist one signal-object possessing an ongoing
identity and living through each intermediate spacetime point described
by the solution.
\end{quote}

\subsection*{Antenna}

At a receiving location, an antenna may undergo a local electrical
response.

That response is a new registration.

It is not necessary to describe the event as the capture of an enduring
object that has survived an intermediate journey.

\subsection*{Filtering and tuning}

A tuned circuit or digital filter preferentially responds to selected
features of the local input.

Tuning is therefore a property of the receiver's local response.

It need not be interpreted as selecting one traveler from a collection
of travelers moving through space.

\subsection*{Demodulation}

Demodulation transforms one local electrical pattern into another.

Amplitude, frequency, phase, or digitally encoded structure can thereby
be converted into baseband electrical information.

Again, the operation is local.

\subsection*{Output}

Amplifiers, speakers, displays, memory devices, and control systems
produce further local state changes.

The complete device-side process can therefore be represented as

\[
\text{local antenna response}
\rightarrow
\text{local selection}
\rightarrow
\text{local demodulation}
\rightarrow
\text{local output}.
\]

% ============================================================
\section{A Diagram of the Distinction}
% ============================================================

\begin{figure}[ht]
\centering

\begin{tikzpicture}[
node distance=0.85cm and 0.65cm,
box/.style={
draw,
rounded corners,
align=center,
minimum height=1.05cm,
minimum width=2.55cm,
fill=gray!8
},
arrow/.style={-{Latex[length=2.2mm]},thick}
]

\node[box] (source) {
Transmitter-side\\
registration
};

\node[box,right=of source] (law) {
Maxwellian / relativistic\\
lawful relation
};

\node[box,right=of law] (receive) {
Receiver-side\\
registration
};

\node[box,right=of receive] (process) {
Local processing\\
and output
};

\draw[arrow] (source) -- (law);
\draw[arrow] (law) -- (receive);
\draw[arrow] (receive) -- (process);

\begin{scope}[on background layer]
\node[
draw,
dashed,
rounded corners,
fit=(source)(law)(receive)(process),
inner sep=9pt,
label=below:{
\small
The arrows denote lawful and engineering relations, not an occupied
biography of a persisting carrier.
}
] {};
\end{scope}

\end{tikzpicture}

\caption{
Radio reception under the no-intermediate-life reading.
The source--receiver relation remains fully calculable.
The diagram does not assign a persisting signal-object to intermediate
locations.
}

\label{fig:radiochain}

\end{figure}

% ============================================================
\section{Finite Delay Without Intermediate Life}
% ============================================================

The refusal of intermediate life does not imply instantaneous reception.

Source-side and receiver-side registrations occupy different spacetime
locations.

The causal relation between them is constrained by relativistic
structure.

The receiver cannot register a source-related outcome arbitrarily early.

The relevant distance, geometry, and causal structure remain part of the
physical description.

\begin{proposition}[Delay does not imply biography]
\label{prop:delay}

A nonzero temporal separation between source-side and receiver-side
registrations does not by itself establish that a persisting
signal-object possesses an intermediate life during that interval.

\end{proposition}

